% Exercises for Chapter 8 (multinomial and Gamma laws; transformations of random variables).
% Statements only: solutions are kept privately. Throughout, (Omega, F, P) is a probability space,
% Gamma(alpha, lambda) has density lambda^alpha x^(alpha-1) e^(-lambda x) / Gamma(alpha) on (0, inf)
% (lambda is a rate), and "seed s" means numpy.random.default_rng(s).

\begin{exercise}[The multinomial law]\label{ex:tr-multinomial}
\medium (in the style of Ross, Ch.~6) Let
$\underline X=(X_1,\dots,X_N)\sim\mathrm{Multinomial}(n,p_1,\dots,p_N)$, with $p_{N+1}=1-\sum_ip_i$
and $X_{N+1}=n-\sum_iX_i$. You may use \noteref{eq:tr-multinomial},
\noteref{rem:tr-multinomial-marginals}, and that $B(n,p)$ has variance $np(1-p)$.
\begin{enumerate}[label=(\alph*)]
  \item Complete the computation that the notes leave to the reader in \noteref{exa:tr-urn}, b):
    compute $\mathbb P(\underline X=(5,3,k))$ for $k=0,1,2$, and check that their sum is the value
    $0.0567$ found with the merged vector $\underline X'$. Explain why merging two types of outcome
    gives again a multinomial vector, so that the two methods must agree.
  \item For $i\ne j$ in $\{1,\dots,N\}$, show that $X_i+X_j\sim B(n,p_i+p_j)$, and deduce that
    $\operatorname{cov}(X_i,X_j)=-np_ip_j$. Why is a negative sign to be expected?
  \item Let $0<p_1<1$ and $0\le k\le n$. Show that, given $X_1=k$, the law of $X_2$ is
    $B\big(n-k,\frac{p_2}{1-p_1}\big)$, that is,
    $\mathbb P(X_2=j\mid X_1=k)=\binom{n-k}j\big(\frac{p_2}{1-p_1}\big)^j\big(1-\frac{p_2}{1-p_1}\big)^{n-k-j}$.
    In the urn of \noteref{exa:tr-urn}, compute $\mathbb P(X_2=3\mid X_1=5)$ in this way and as a
    ratio of two answers of the worked exercise.
  \item A fair die is rolled $12$ times. Find the probability that each face appears exactly
    twice, and the probability that more sixes than ones appear. \emph{Hint for the second:}
    symmetry, and the trinomial law of (number of ones, number of sixes). Check both answers by
    simulating $10^6$ series of $12$ rolls (seed $2026$).
\end{enumerate}
\end{exercise}

\begin{exercise}[The annulus, explicitly]\label{ex:tr-annulus}
\easy Let $(X,Y)$ be uniform on $D=\{(x,y):1\le x^2+y^2\le9\}$, as in \noteref{exa:tr-annulus},
whose explicit computations are left to the reader. You may use \noteref{thm:pm-tonelli},
\noteref{thm:pm-fubini}, the marginal density found in the notes, and, in (c), polar coordinates,
$\iint_Dh(x,y)\,\dd x\,\dd y=\int_1^3\int_0^{2\pi}h(r\cos\theta,r\sin\theta)\,r\,\dd\theta\,\dd r$
for $h\ge0$ or integrable on $D$.
\begin{enumerate}[label=(\alph*)]
  \item Compute $\mathbb E[X]$, $\mathbb E[Y]$ and $\mathbb E[XY]$ as explicit integrals, saying
    where you use Tonelli's and where Fubini's theorem.
  \item Show that $\mathbb E[X^2]=\frac52$, using $f_X$ and $\int_{-r}^rx^2\sqrt{r^2-x^2}\,\dd x=\frac{\pi r^4}8$.
  \item Show that $\mathbb E[X^2Y^2]=\frac{91}{24}$ and $\operatorname{cov}(X^2,Y^2)=-\frac{59}{24}$.
    Why does this prove again that $X$ and $Y$ are not independent?
  \item Find $a,b\in\mathbb R$ with $\mathbb P(X>a,\,Y>b)\ne\mathbb P(X>a)\,\mathbb P(Y>b)$,
    without computing any integral.
\end{enumerate}
\end{exercise}

\begin{exercise}[Poisson conditionals, dependent]\label{ex:tr-poisson-conditionals}
\medium This exercise continues \noteref{exa:tr-poisson-joint}, now with $0<\alpha<1$:
$f_{(X,Y)}(n,m)=c\,\frac{\nu^n\mu^m}{n!\,m!}\alpha^{nm}$ on $\mathbb N^2$.
\begin{enumerate}[label=(\alph*)]
  \item Show that $X$ and $Y$ are not independent, by comparing $f_{(X,Y)}(0,0)f_{(X,Y)}(1,1)$ with
    $f_{(X,Y)}(0,1)f_{(X,Y)}(1,0)$. Why are these two products equal for a joint density of
    independent random variables?
  \item Show that $X$ and $Y$ have finite moments of every order. Using b) of the worked
    exercise, show that $\sum_nn\,\mathbb P(X=n\mid Y=m)=\nu\alpha^m$, and deduce that
    $\mathbb E[X]=\mathbb E[\nu\alpha^Y]$, $\mathbb E[XY]=\mathbb E[Y\nu\alpha^Y]$, so that
    $\operatorname{cov}(X,Y)=\operatorname{cov}(Y,\nu\alpha^Y)$.
  \item Show that $\operatorname{cov}(X,Y)<0$. \emph{Hint:} for functions $\varphi,\psi$ with
    $\varphi(Y),\psi(Y),\varphi(Y)\psi(Y)\in L^1$,
    $\operatorname{cov}(\varphi(Y),\psi(Y))=\frac12\sum_{m,m'}f_Y(m)f_Y(m')\big(\varphi(m)-\varphi(m')\big)\big(\psi(m)-\psi(m')\big)$.
  \item This part is computational. For $\nu=\mu=2$ and $\alpha\in\{1,0.75,0.5,0.25\}$, compute
    $c$ (summing the series for $c^{-1}$ of the worked exercise), $\mathbb E[X]$,
    $\operatorname{var}(X)$ and the correlation coefficient $\rho(X,Y)$. Check that
    $c=e^{-4}$ for $\alpha=1$. How does $\rho$ depend on $\alpha$?
\end{enumerate}
\end{exercise}

\begin{exercise}[Moments and scale of the Gamma law]\label{ex:tr-gamma-scale}
\easy Let $X\sim\Gamma(\alpha,\lambda)$. You may use \noteref{prop:tr-gamma-function} and
\noteref{thm:tr-density}.
\begin{enumerate}[label=(\alph*)]
  \item For real $k>-\alpha$, show that
    $\mathbb E[X^k]=\frac{\Gamma(\alpha+k)}{\Gamma(\alpha)\lambda^k}$, and that $\mathbb E[X^k]=+\infty$ for
    $k\le-\alpha$. Recover \noteref{prop:tr-gamma-moments}, and show that
    $\mathbb E[1/X]=\frac\lambda{\alpha-1}$ for $\alpha>1$.
  \item Let $c>0$. Show that $cX\sim\Gamma(\alpha,\frac\lambda c)$. Deduce that if
    $Z\sim\Gamma(\alpha,1)$, then $Z/\lambda\sim\Gamma(\alpha,\lambda)$: $\frac1\lambda$ is a scale
    parameter and $\lambda$ a rate (\noteref{rem:tr-gamma-special}).
  \item For $\alpha>1$, show that $f_X$ has a unique maximum, at $\frac{\alpha-1}\lambda$. Compute
    the limits $f_X(0^+)$ and $f_X'(0^+)=\lim_{x\to0^+}f_X'(x)$ for $0<\alpha<1$, $\alpha=1$,
    $1<\alpha<2$, $\alpha=2$ and $\alpha>2$, and compare with \noteref{fig:tr-gamma} and
    \noteref{rem:tr-gamma-special}.
  \item Using $\Gamma(\frac12)=\sqrt\pi$ (\noteref{exa:tr-chi-square-end}), show by induction that
    $\Gamma(n+\frac12)=\frac{(2n)!}{4^nn!}\sqrt\pi$ for $n\in\mathbb N$, and compute
    $\mathbb E[\sqrt X]$ for $X\sim\mathrm{Exp}(1)$.
\end{enumerate}
\end{exercise}

\begin{exercise}[Wear and rejuvenation]\label{ex:tr-ageing}
\hard This exercise proves \noteref{rem:tr-ageing}, which the notes state without proof. Let
$X\sim\Gamma(\alpha,\lambda)$, $S(x)=\mathbb P(X>x)$, and let $h(x)=f_X(x)/S(x)$ for $x>0$ be the
\emph{hazard rate} of $X$.
\begin{enumerate}[label=(\alph*)]
  \item Show that $S$ is continuous on $[0,+\infty)$ and of class $C^1$ on $(0,+\infty)$, with
    $S'=-f_X$, and that $h$ is integrable on $(0,t)$ for every $t>0$. Deduce that
    \[
      \mathbb P(X>s+t\mid X>s)=\exp\Big(-\int_s^{s+t}h(u)\,\dd u\Big)\qquad(s\ge0,\ t>0).
    \]
  \item Show that, for $x>0$,
    $\frac1{h(x)}=\int_0^{+\infty}\big(1+\frac ux\big)^{\alpha-1}e^{-\lambda u}\,\dd u$. Deduce that
    $h\equiv\lambda$ if $\alpha=1$, that $h$ is strictly increasing if $\alpha>1$ and strictly
    decreasing if $\alpha<1$, and that $h(x)\to\lambda$ as $x\to+\infty$ in every case.
  \item Deduce the three statements of \noteref{rem:tr-ageing}, for all $s,t>0$. Why do they fail
    for $s=0$?
  \item Compute $\mathbb P(X>s+1\mid X>s)$ for $\lambda=1$, $s\in\{0,1,5\}$ and
    $\alpha\in\{\frac12,1,3\}$. What is the limit as $s\to+\infty$, and why?
\end{enumerate}
\end{exercise}

\begin{exercise}[Simulating a Poisson process]\label{ex:tr-poisson-process}
\medium This is a computational exercise. Let $U_1,U_2,\dots$ be i.i.d.\ $\mathrm{Exp}(\lambda)$,
$T_n=U_1+\dots+U_n$ and $N_t=\#\{n\ge1:T_n\le t\}$ for $t\ge0$, as in
\noteref{rem:tr-poisson-process}. You may use that $T_n\sim\Gamma(n,\lambda)$
(\noteref{prop:rv-notable-sums}).
\begin{enumerate}[label=(\alph*)]
  \item Show that $\{N_t\ge n\}=\{T_n\le t\}$, and deduce from \noteref{prop:tr-erlang} that
    $N_t\sim\mathcal P(\lambda t)$. In particular $N_t<+\infty$ a.s.
  \item Let $t>0$. Show that $\mathbb P(T_1\le s\mid N_t=1)=\frac st$ for $0\le s\le t$: given that
    exactly one arrival occurs in $[0,t]$, its time is uniform on $[0,t]$. \emph{Hint:}
    $\{T_1\le s,\,N_t=1\}=\{U_1\le s,\,U_1+U_2>t\}$, and $(U_1,U_2)$ has the density
    $\lambda^2e^{-\lambda(u_1+u_2)}$ on $(0,+\infty)^2$ (\noteref{thm:pm-density-independence}).
  \item Take $\lambda=2$ and seed $2026$. Simulate $10^5$ paths and compute the
    Kolmogorov--Smirnov distance $\sup_t|F^{(n)}(t)-F_5(t)|$ between the empirical distribution
    function of the $10^5$ values of $T_5$ and the Erlang distribution function $F_5$. Compare it
    with $1.36/\sqrt n$, which a correct model exceeds with probability about $5\%$ for large $n$.
    With the same generator, simulate $10^5$ values of $N_3$ and compute the largest difference
    between their relative frequencies and the density of $\mathcal P(6)$.
  \item How many paths $n$ are needed so that the estimate $\hat p$ of $p=\mathbb P(T_5\le2)$ has an
    asymptotic $95\%$ confidence half-width $1.96\sqrt{\hat p(1-\hat p)/n}$ of at most $0.005$,
    whatever the value of $p$? Run this many paths (seed $2027$) and compare $\hat p$ with the exact
    value given by \noteref{prop:tr-erlang}.
  \item Check (b) by simulation for $t=1$ (seed $2028$, $10^5$ paths): estimate $\mathbb P(N_1=1)$,
    and compute the Kolmogorov--Smirnov distance between the values of $T_1$ on the paths with
    $N_1=1$ and the uniform law on $[0,1]$.
\end{enumerate}
\end{exercise}

\begin{exercise}[Affine maps of a Gaussian; chi-square by simulation]\label{ex:tr-chi-square}
\easy Part (a) is the computation left to the reader in \noteref{prop:tr-affine-gaussian}; parts (b) and
(c) are computational.
\begin{enumerate}[label=(\alph*)]
  \item Do the direct computation that the notes leave as an exercise in
    \noteref{prop:tr-affine-gaussian}: for an absolutely continuous $X$ and $Y=\sigma X+\mu$,
    compute $F_Y$ in terms of $F_X$ for $\sigma>0$ and for $\sigma<0$, and recover
    $f_Y(y)=\frac1{|\sigma|}f_X\big(\frac{y-\mu}\sigma\big)$ at the points where $f_X$ is continuous.
  \item Let $Z\sim\mathcal N(0,1)$. Simulate $10^6$ values of $Z^2$ (seed $2026$) and compute their
    Kolmogorov--Smirnov distance to the distribution function of $\Gamma(\frac12,\frac12)$; compare
    with $1.36/\sqrt n$. Make a histogram on $[0,4]$ with $40$ bins, and compare the height of the
    first bin with the density at the midpoint $0.05$ and with the bin average
    $\frac{F(0.1)-F(0)}{0.1}$. Which of the two should the height estimate, and why do they differ so
    much here?
  \item $\Gamma(\frac12)=\mathbb E[V^{-1/2}]$ for $V\sim\mathrm{Exp}(1)$. Show that this estimator
    of $\sqrt\pi$ has infinite variance. Show that also $\Gamma(\frac12)=2\,\mathbb E[\sqrt V]$, and
    that $\operatorname{var}(2\sqrt V)=4-\pi$. How many samples give an asymptotic $95\%$
    half-width of $0.002$ for the second estimator? Compute both estimators with $10^3,\dots,10^6$
    samples of $V$ (seed $2026$) and comment.
\end{enumerate}
\end{exercise}

\begin{exercise}[Piecewise monotone transformations]\label{ex:tr-piecewise}
\medium This is a computational exercise. You may use \noteref{cor:tr-piecewise}.
\begin{enumerate}[label=(\alph*)]
  \item Let $X$ be uniform on $(-1,2)$ and $Y=X^2$. With the pieces $I_1=(-1,0)$ and $I_2=(0,2)$,
    show that $f_Y(y)=\frac1{3\sqrt y}$ for $0<y<1$, $f_Y(y)=\frac1{6\sqrt y}$ for $1<y<4$, and
    $f_Y(y)=0$ otherwise. Check this by computing $F_Y$ directly, and explain the jump of $f_Y$ at
    $y=1$ (\noteref{fig:tr-transform}).
  \item Let $X$ be uniform on $(0,2\pi)$ and $Y=\sin X$. Using the three pieces
    $(0,\frac\pi2)$, $(\frac\pi2,\frac{3\pi}2)$, $(\frac{3\pi}2,2\pi)$, show that
    $f_Y(y)=\frac1{\pi\sqrt{1-y^2}}$ for $-1<y<1$ (the \emph{arcsine law}). Which pieces contribute
    for $y>0$, and which for $y<0$? Find $F_Y$.
  \item Write a function that evaluates the formula of the corollary from $f_X$, $g$ and a list of
    pieces on which $g$ is strictly monotone, inverting $g$ on each piece numerically by bisection
    and differentiating the inverse as $1/g'(g_i^{-1}(y))$. Check it against the closed forms of (a)
    and (b), and compare (a) and (b) with $10^6$ simulated values (seeds $2026$ and $2027$) through
    the Kolmogorov--Smirnov distance.
\end{enumerate}
\end{exercise}

\begin{exercise}[A lamp and a wall]\label{ex:tr-cauchy}
\medium (in the style of Ross, Ch.~5) A lamp stands at distance $1$ from an infinite straight wall.
It is pointed at a random angle $\Theta$, uniform on $(-\frac\pi2,\frac\pi2)$ and measured from the
perpendicular to the wall, so that the light hits the wall at the signed distance $Y=\tan\Theta$
from the foot of the perpendicular.
\begin{enumerate}[label=(\alph*)]
  \item Show that $Y$ has the density $\frac1{\pi(1+y^2)}$ on $\mathbb R$, the \emph{Cauchy law}.
  \item Show that $\mathbb E[|Y|]=+\infty$, so that $\mathbb E[Y]$ is not defined, although the
    density is symmetric.
  \item The lamp is replaced by a double beam along a line through the lamp, at an angle $\Theta$
    uniform on $(-\pi,\pi)$; the line hits the wall at $\tan\Theta$ (for $\Theta\ne\pm\frac\pi2$).
    Use \noteref{cor:tr-piecewise} with three pieces to show that $\tan\Theta$ is again Cauchy, and
    find the image $g(I_i)$ of each piece.
  \item Show that $1/Y$ is also Cauchy.
  \item Simulate $10^6$ values of $\tan\Theta$, $\Theta$ uniform on $(-\pi,\pi)$ (seed $2026$),
    and print the mean and the median of the first $10^2,10^3,\dots,10^6$ values. Explain what you see.
\end{enumerate}
\end{exercise}

\begin{exercise}[Calls at a switchboard]\label{ex:tr-arrivals}
\medium (in the style of Ross, Ch.~5) Calls arrive at a switchboard as a Poisson process of rate
$\lambda=6$ per hour (\noteref{rem:tr-poisson-process}): the times between successive calls are
i.i.d.\ $\mathrm{Exp}(6)$, in hours, and the $n$-th call arrives at $T_n\sim\Gamma(n,6)$.
\begin{enumerate}[label=(\alph*)]
  \item What is the probability that the tenth call arrives within the first hour?
  \item Find the mean and the standard deviation of the arrival time of the tenth call, in minutes.
  \item What is the probability that no call arrives in the first $20$ minutes? Given that none
    arrived in the first $20$ minutes, what is the probability that none arrives in the next $20$
    minutes? Which property of the exponential law is this?
  \item Find, numerically, the time $t$ by which the tenth call has arrived with probability
    $0.9$, by solving $F_{10}(t)=0.9$ with \noteref{prop:tr-erlang}.
  \item A second switchboard, independent of the first, receives calls as a Poisson process of
    rate $\mu=4$ per hour; its second call arrives at a time $B\sim\Gamma(2,\mu)$. Show that the
    probability that the first call at the first switchboard arrives before the second call at the
    second one is $1-\frac{\mu^2}{(\lambda+\mu)^2}=0.84$. \emph{Hint:} condition on $B$, that is, use
    \noteref{thm:pm-independence-product} and \noteref{thm:pm-tonelli} to write
    $\mathbb P(T_1<B)=\int_0^{+\infty}\mathbb P(T_1<t)\,f_B(t)\,\dd t$. Check by simulation
    ($10^6$ pairs, seed $2026$).
\end{enumerate}
\end{exercise}
